\scanpage{193}\mspage{142}
elles définissent des objets isomorphes de $\mathrm{HOT}_W(\Delta^1)$.

\textcircled{d} Plus précisément \unsure{encore|même}, les quatre diagrammes dans $\mathrm{Cat}$, de type $\Gamma^\circ$, où

\diagram{\(\Gamma={}\begin{tikzcd}[ampersand replacement=\&,row sep=.4em,column sep=1.2em]
\& 1\arrow[dl,"\alpha_1"'] \& \& \\
0 \& \& 3\arrow[ul,"i'_1"']\arrow[dl,"i'_2"] \& 4\arrow[l,"\varphi"']\\
\& 2\arrow[ul,"\alpha_2"] \& \&
\end{tikzcd},\quad\alpha_1i'_1\varphi=\alpha_2i'_2\varphi,\)}

dans lesquels s'insèrent les quatre flèches qu'on vient d'envisager, savoir

\diagram{\begin{tikzcd}[ampersand replacement=\&,row sep=.3em,column sep=1.4em]
\& \& X_1\arrow[dr,"i_1"] \& \& \\
\widehat Q={} \& X_0\arrow[ur,"\alpha_1"]\arrow[dr,dash,"\alpha_2"']\arrow[rr,dotted] \& \& \unsure{\mathscr X|\mathcal X}_Q\arrow[r,"\varphi_Q"] \& Y\\
\& \& X_2\arrow[ur,"i_2"'] \& \&
\end{tikzcd}}

% slanted
\margin{\uline{NB} On devrait \deleted{\unsure{et}} encore compléter ces diagrammes \deleted{par les} \unsure{(bien comm.)} par les flèches en pointillés\ldots}

\diagram{\begin{tikzcd}[ampersand replacement=\&,row sep=.3em,column sep=1.4em]
\& \& X_1^\circ\arrow[dr,"i'_1"] \& \& \\
\widehat{Q^\circ}_{\deleted{\unsure{\circ}}}={} \& X_0^\circ\arrow[ur,"\alpha_1^\circ"]\arrow[dr,"\alpha_2^\circ"']\arrow[rr,dotted] \& \& \mathscr X_{Q^\circ}\arrow[r] \& Y^\circ\\
\& \& X_2^\circ\arrow[ur,"i'_2"'] \& \&
\end{tikzcd}}

\diagram{\begin{tikzcd}[ampersand replacement=\&,row sep=.3em,column sep=1.4em]
\& \& X_1\arrow[dr] \& \& \\
\widehat Q^{\sim}={} \& X_0\arrow[ur]\arrow[dr]\arrow[rr,dotted] \& \& \widetilde{\mathscr X}_{\unsure{\overline Q|Q}}\arrow[r] \& Y\\
\& \& X_2\arrow[ur] \& \&
\end{tikzcd}}

\diagram{\begin{tikzcd}[ampersand replacement=\&,row sep=.3em,column sep=1.4em]
\& \& X_1^\circ\arrow[dr,Rightarrow] \& \& \\
\widehat{Q^\circ}^{\sim}={} \& X_0^\circ\arrow[ur]\arrow[dr]\arrow[rr,dotted] \& \& \widetilde{\mathscr X}_{Q^\circ}\arrow[r] \& Y^\circ\\
\& \& X_2^\circ\arrow[ur,dash] \& \&
\end{tikzcd}}

\deleted{\unsure{déjà|déf.}} sont isomorphes dans $\mathrm{Hot}_W(\Gamma)$. De façon plus précise encore, on a les isomorphismes dans \uline{$\mathrm{Hom}$}$(\Gamma^\circ,\mathrm{Cat})$ :

\[
\begin{aligned}
(\widehat Q^{\sim})^\circ &\simeq \widehat{Q^\circ}\\
(\widehat{Q^\circ}^{\sim})^\circ &\simeq \widehat Q,
\end{aligned}
\]

\margin{d'où en l'appliquant à $Q^\circ$ au lieu de $Q$}

ce qui implique (vu les isom. canoniques valables dans $\mathrm{HOT}_W(\unsure{S|D})$ pour \deleted{toute catégorie} objet $S$ \deleted{dans} $\mathrm{Cat}$ et tout $X_*\in\mathrm{Ob}\underline{\mathrm{Hom}}(S,\mathrm{Cat})$, $X_*\simeq X_*^\circ$${}^{\star}$) qu'on a les isom. canoniques
% slanted
\margin{via le foncteur $\unsure{\Sigma|Z}$ utilisé plus haut\ldots}

\[
\left.\begin{aligned}
\widehat Q&\simeq\widehat{Q^\circ}^{\sim}\\
\widehat{Q^\circ}&\simeq\widehat Q^{\sim}
\end{aligned}\right|\quad\text{dans }\mathrm{HOT}_W(\Gamma)
\]

mais aussi directement

\[
\widehat Q\simeq\widehat{Q^\circ},\quad \widehat Q^{\sim}\simeq\widehat{Q^\circ}^{\sim}\quad\text{dans }\mathrm{HOT}_W(\Gamma)
\]
% archive sheet number: 193
